%----------------------------------------------------------------------------------------
%	                        MAIN 
%----------------------------------------------------------------------------------------

%----------------------------------------------------------------------------------------
%	CONFIGURATIONS
%----------------------------------------------------------------------------------------
\documentclass[10pt, a4paper]{article}
%----------------------------------------------------------------------------------------
%	MERLIN CAMBERLIN ! ALL RIGHTS RESERVED
%----------------------------------------------------------------------------------------

%----------------------------------------------------------------------------------------
%	PACKAGES ET AUTRES CONFIGURATIONS
%----------------------------------------------------------------------------------------
\usepackage[french]{babel} 
\usepackage[utf8]{inputenc}
%\usepackage[latin1]{inputenc}
\usepackage[T1]{fontenc}

\usepackage{supertabular}% pour séparer un tableau lorsqu'il est sur plusieurs pages
\usepackage{ulem} %pour souligner
\usepackage[bottom]{footmisc} % pour mettre les notes de bas de page bien en bas
%%% Pour les couleurs
\usepackage{pdflscape}
\usepackage[svgnames]{xcolor}
\usepackage{colortbl}%
    \newcommand{\myrowcolour}{\rowcolor[gray]{0.925}}
\usepackage{booktabs}
\definecolor{lightgray}{gray}{0.5}

%Pour supprimer l'indentation automatique
\usepackage{parskip}
\usepackage{comment} % Pour ajouter des commentaires sur plusieurs lignes avec \begin{comment} ... \end{comment}



%%% Pour la mise en page par défaut
\usepackage{fullpage}
\usepackage[left=2cm,right=2cm,top=2cm,bottom=2cm]{geometry}
\usepackage{lastpage}%Donne la dernière page

\usepackage{setspace} % gérer l'interligne

%%% Pour écrire du texte sur plusieurs colonnes
\usepackage{multicol}
\usepackage{enumitem} % why not, item
\frenchbsetup{StandardLists=true} % à inclure si on utilise \usepackage[french]{babel} pour des bubulles

%%% Pour les circuits électriques
\usepackage{tikz}
\usepackage{circuitikz}


%%% Pour créer un en-tête et un pied de page
\usepackage{fancyhdr} % Needed to define custom headers/footers
\pagestyle{fancy} % Enables the custom headers/footers

\renewcommand{\headrulewidth}{0.pt} % No header rule
\renewcommand{\footrulewidth}{0.pt} % Thin footer rule

\renewcommand{\sectionmark}[1]{\markboth{#1}{}} % Removes the section number from the header when \leftmark is used

%\nouppercase\leftmark % Add this to one of the lines below if you want a section title in the header/footer

% Headers
\lhead{} % Left header
\chead{} % Center header - currently printing the article title
\rhead{} % Right header

% Footers
%\lfoot{\includegraphics[scale=0.18]{Images/exemple.jpg}}
%\cfoot{\color{lightgray}{\nouppercase\leftmark }% Add this to one of the lines below if you want a section title in the header/footer} % Center footer
\fancyfoot[CE,CO]{\color{lightgray} \nouppercase \leftmark}
\fancyfoot[LE,RO]{\thepage}
%\rfoot{\footnotesize Page \thepage\ sur \pageref{LastPage}} % Right footer, "Page 1 of 2"
\usepackage{etoolbox}
\patchcmd{\chapter}{\thispagestyle{plain}}{\thispagestyle{fancy}}{}{}


%%% Bibliographie
\usepackage[sorting=none]{biblatex} 
    \addbibresource{bibliographie.bib}
\usepackage[autostyle=true]{csquotes} 

%%% Pour insérer des mathématiques
\usepackage{amsmath,amsfonts,amsthm,amssymb}
\usepackage{eqnarray}
\usepackage{systeme} % pour faire des systèmes d'équations
\usepackage{numprint} % Pour que les chiffres s'écrivent seul en notation scientifique


%%% Pour insérer des caractères spéciaux
\usepackage{gensymb}    % Pour insérer le signe degré
\usepackage{eurosym}    % Pour insérer le signe €

%%% Pour insérer des images
\usepackage{graphicx}
\usepackage{float}          % Pour obliger une position d'une figure \begin{figure}[H]
\usepackage[tt]{titlepic}   % Package pour insérer une image dans la première page
\usepackage{subfigure}  % Permet d'avoir des sous-images
\usepackage{wrapfig} % Permet d'avoir du texte autour d'une image

%%% Pour insérer un pdf
\usepackage{pdfpages}
%%% Pour insérer des URL
\usepackage{url}

%%% Package pour insérer du code 
\usepackage{listings}
\usepackage{color}
 
\definecolor{codegreen}{rgb}{0,0.6,0}
\definecolor{codegray}{rgb}{0.5,0.5,0.5}
\definecolor{codepurple}{rgb}{0.58,0,0.82}
\definecolor{backcolour}{rgb}{0.95,0.95,0.92}
 
\lstdefinestyle{styleMerlin}{
    language = C,% ou java
    backgroundcolor=\color{backcolour},  
    commentstyle=\color{codegreen},%\color[rgb]{0.133,0.545,0.133},
    %identifierstyle=\ttfamily,
    keywordstyle=\color{magenta},%\color[rgb]{0,0,1},
    numberstyle=\tiny\color{codegray},
    stringstyle=\color{codepurple},%\color[rgb]{0.627,0.126,0.941},
    basicstyle=\footnotesize,%\scriptsize
    breakatwhitespace=false,         
    breaklines=true,                 
    captionpos=b,                    
    keepspaces=true,                 
    numbers=left,                    
    numbersep=5pt,                  
    showspaces=false,                
    showstringspaces=false,
    showtabs=false,                  
    tabsize=2 % ajouter la virgule ici et l'enlever à la derniere ligne
    %upquote=true,
    %aboveskip={1.5\baselineskip},
    %columns=fullflexible,
    %extendedchars=true,
}
\lstset{style=styleMerlin}
%\begin{lstlisting}[language=C, caption=Python example]
%#include <stdio.h>
% ...
%\end{lstlisting}
%----------------------------------------------------------------------------------------
%	MERLIN CAMBERLIN ! ALL RIGHTS RESERVED
%----------------------------------------------------------------------------------------
\usepackage{geometry}
\geometry{hmargin= 2cm,vmargin=2cm}
\bibliography{bibliographie}
\newtheorem{theorem}{Formule}
\renewcommand{\arraystretch}{1.8}
\setlength\fboxrule{0.1cm}
\setlength\fboxsep{0.5cm}
\addtocounter{tocdepth}{1} % afficher les subsubsection dans la table des matières
\setcounter{secnumdepth}{4} % numéroter les subsubsection
\setcounter{tocdepth}{4} % numéroter les subsubsection
\newcommand{\te}{\texttt}

%----------------------------------------------------------------------------------------
%	DOCUMENT
%----------------------------------------------------------------------------------------
\begin{document}

%----------------------------------------------------------------------------------------
%	PAGE DE GARDE
%----------------------------------------------------------------------------------------

%----------------------------------------------------------------------------------------
%	PAGE DE GARDE
%----------------------------------------------------------------------------------------
\begin{titlepage}

  \begin{sffamily}
  
  \begin{center}  
        
        \begin{figure}[h]
            \begin{minipage}[c]{0.5\textwidth}
            \flushleft
            \includegraphics[scale=0.8]{Structure/LogoEPL.jpg}
            \end{minipage}%
            \begin{minipage}[c]{0.5\textwidth}
            \includegraphics[scale=0.37]{Structure/LogoUCLouvain2.jpg}
            \flushright 
            \end{minipage}
        \end{figure}
        
\vfill

        %%% TITRE
        \textsc{\LARGE LINGI 1131  - Computer language concepts -  2020}\\[0.3cm]
        \rule{16cm}{1pt}
        \medskip
        \\
        { \huge \bfseries Captain Sonoz Project \\[0.2cm] }
        \rule{16cm}{1pt}
        \\
    
        \begin{figure}[h!]
        \centering
        \includegraphics[scale = 0.8]{Images/couverture2.jpg}
        \end{figure}
    
\vfill
\large   
    
    %%% NOM & Prénom & NOMA \\ 
        \begin{table}[ht!]
\centering
\large
            \begin{tabular}{lcr}
                     & \Large{Auteurs} & \\
                \LARGE{\textsc{Camberlin}} & \Large{Merlin} & \Large{0944-17-00} \\
                \LARGE{\textsc{Pisvin}} & \Large{Arthur} & \Large{-17-00} \\

            \end{tabular}
  
        \end{table}
\vspace{3cm}
    \begin{center}
        Soumis le : \today
    \end{center}
    
    
\vfill

%Professeur et tuteur
\begin{tabular}{p{.5\textwidth}p{.5\textwidth}}
    \flushleft Professeur : \\
    & \flushright P. Van Roy\\
 
\end{tabular} 
    
\end{center}
\end{sffamily}
\end{titlepage}


%----------------------------------------------------------------------------------------
%	TABLE DES MATIÈRES
%----------------------------------------------------------------------------------------
%\tableofcontents
\pagenumbering{arabic}
\setcounter{page}{1}

%----------------------------------------------------------------------------------------
%	CONTENU
%----------------------------------------------------------------------------------------
\section{Stratégie de chaque joueur}

\section{Conception et  choix d'implémentation}
Le record \texttt{state} a été utilisé dans tout le code pour représenter l'état courant d'un sous-marin:
\begin{itemize}
    \item state ( 
                \begin{itemize}
                    \item id( id<idNum> color:<color> name:Name) 
                    \item position: <position>
                    \item lastPosition: <position>
                    \item direction: <direction> 
                    \item surface: <true> | <false>
                    \item loads: load( mine: x missile: y drone: z sonar: u) 
                    \item weapons: weapons( mine: x missile: y drone: z sonar:u)
                    \\)
                \end{itemize}
\end{itemize}

\section{Extension(s) implémentée(s)}


\end{document}
